\begin{textsample}{POS Dim 8 – human – Score 2.00 – mg/mg\_ef\_er\_7\_p4.txt}  \label{ex:f8_pos_017}
Deve -se observar que o sujeito principal desse processo é o educando e não o currículo ou o professor.
No componente curricular do Ensino Religioso é muito importante que haja um profundo respeito às crenças dos educandos e de suas famílias, de modo que não haja espaço para nenhum tipo de desrespeito e também de proselitismo religioso.
O professor deve levar em \textbf{conta} a situação do processo de desenvolvimento dos educandos, a \textbf{realidade} local e \textbf{regional} ao trabalhar com as unidades temáticas, os objetos de conhecimento e suas habilidades.

% matched lemmas: conta, realidade, regional
\end{textsample}
